\section*{Capítulo \ref{ch:b2:alkene-redox} --- Oxidação e redução dos alcenos}

\begin{solution}{exo:b2:alkene-redox:1}
Metilciclo-hexano; 1-metilciclo-hexan-1-ol; trans-2-metilciclo-hexan-1-ol (adição
syn de H e de B; logo o \ce{OH} e o novo H do mesmo lado, a metila em trans em
relação ao \ce{OH}); óxido de 1-metilciclo-hexeno; cis-1-metilciclo-hexano-1,2-diol;
6-oxo-heptanal \ce{CH3CO(CH2)4CHO}.
\end{solution}

\begin{solution}{exo:b2:alkene-redox:2}
Hidroboração--oxidação: 2-metilpropan-1-ol. Hidratação catalisada por ácido:
2-metilpropan-2-ol.
\end{solution}

\begin{solution}{exo:b2:alkene-redox:3}
trans-2,3-Dimetiloxirano, com as duas metilas em lados opostos do anel. Ele é
quiral (sem plano de simetria) e se forma como mistura racêmica.
\end{solution}

\begin{solution}{exo:b2:alkene-redox:4}
Tratamento redutor: duas moléculas de propanal. Tratamento oxidante: duas
moléculas de ácido propanoico.
\end{solution}

\begin{solution}{exo:b2:alkene-redox:5}
(a) meso-butano-2,3-diol, inativo por ser aquiral. (b) os dióis
(2\textit{R},3\textit{R}) e (2\textit{S},3\textit{S}) em quantidades iguais, inativos
por serem racêmicos.
\end{solution}

\begin{solution}{exo:b2:alkene-redox:6}
Em meio ácido, a água ataca o carbono terciário; em meio básico, o hidróxido
ataca o \ce{CH2}; nos dois casos, termina-se com um \ce{OH} em cada carbono: o
mesmo 2-metilpropano-1,2-diol. Com o metanol, as duas vias dão éteres
diferentes (como no capítulo).
\end{solution}

\begin{solution}{exo:b2:alkene-redox:7}
Hidrogenação com um equivalente de \ce{H2} sobre um catalisador de paládio
envenenado (Lindlar): a adição syn à ligação tripla para no (\textit{Z})-alceno.
\end{solution}

\begin{solution}{exo:b2:alkene-redox:8}
A propanona (3 C) e o butanal (4 C) se unem por seus carbonos carbonílicos,
dando o 2-metil-hex-2-eno, \ce{(CH3)2C=CH-CH2CH2CH3}.
\end{solution}

\begin{solution}{exo:b2:alkene-redox:9}
\ce{BH3} (ou um borano volumoso), depois \ce{H2O2} e \ce{NaOH}: hexan-1-ol. A
hidratação catalisada por ácido passa pelo carbocátion secundário e dá o
hexan-2-ol (Markovnikov), com subprodutos rearranjados.
\end{solution}

\begin{solution}{exo:b2:alkene-redox:10}
O \omterm{def:b2:alkene-redox:epoxide}{perácido}, um eletrófilo, reage mais depressa com a ligação dupla do anel,
mais substituída e mais rica em elétrons. O borano volumoso reage com o
\ce{C=CH2} terminal, menos impedido.
\end{solution}

\begin{solution}{exo:b2:alkene-redox:11}
trans: \omterm{def:b2:alkene-redox:epoxide}{perácido}, depois água em meio ácido (abertura anti). cis: \ce{OsO4}
catalítico com um reoxidante, ou permanganato diluído a frio (syn).
\end{solution}

\begin{solution}{exo:b2:alkene-redox:12}
\ce{C6H10} tem dois graus de insaturação; um equivalente de \ce{H2} significa uma
\ce{C=C}, logo um anel. Um único dialdeído de seis carbonos vem de um anel
aberto em sua ligação dupla: o ciclo-hexeno.
\end{solution}

\begin{solution}{pb:b2:alkene-redox:1}
\textbf{1.} $10 \times 12.011 + 16 \times 1.008 = \qty{136.24}{g/mol}$.
\textbf{2.} $(2 \times 10 + 2 - 16)/2 = 3$.
\textbf{3.} Um anel e duas ligações duplas.
\textbf{4.} \ce{C10H16 + 2H2 -> C10H20}, o 1-isopropil-4-metilciclo-hexano.
\textbf{5.} $1.00/136.24 = \qty{7.34}{mmol}$ de limoneno, \qty{14.68}{mmol} de \ce{H2}.
\textbf{6.} $V = nRT/p = 0.01468 \times 8.314 \times 298.15/10^5 = \qty{3.64e-4}{m^3}$,
\qty{364}{mL}.
\textbf{7.} Não: os carbonos 1 e 4 do anel têm, cada um, dois caminhos idênticos
pelo anel; os isômeros cis e trans são aquirais.
\textbf{8.} O metanal \ce{HCHO} (1 C), vindo da extremidade \ce{=CH2}, e um único
composto \ce{C9}, o 3-acetil-6-oxo-heptanal \ce{CH3CO-CH2CH2-CH(COCH3)-CH2-CHO}.
\textbf{9.} Seus dois carbonos pertencem ao mesmo anel: romper a ligação dupla
abre o anel sem separar nada.
\textbf{10.} O metanal (um gás).
\textbf{11.} Os grupos aldeído se tornam ácidos carboxílicos (o metanal vai até
\ce{CO2}); as cetonas não mudam.
\textbf{12.} A cetona e o aldeído da cadeia \ce{C9} estavam unidos no anel; a
segunda cetona (\ce{COCH3}) e o metanal eram as duas extremidades da ligação
dupla da cadeia lateral.
\textbf{13.} Um (o C4 do anel); ele não é afetado pela \omterm{def:b2:alkene-redox:cleavage}{ozonólise} e permanece como
estereocentro do produto \ce{C9}.
\textbf{14.} O \ce{C=CH2} exocíclico, menos impedido que a ligação dupla
trissubstituída do anel.
\textbf{15.} 2-(4-metilciclo-hex-3-en-1-il)propan-1-ol: \ce{CH2OH} no antigo
carbono terminal.
\textbf{16.} Primário.
\textbf{17.} O \ce{OH} vai para o carbono menos substituído, ao contrário da
orientação prevista pela regra de Markovnikov para a hidratação.
\textbf{18.} Duas configurações no novo centro, dando dois diastereoisômeros (com
o centro já existente).
\textbf{19.} A hidratação catalisada por ácido (ou outra hidratação de
Markovnikov): o álcool terciário, o $\alpha$-terpineol.
\textbf{20.} A ligação dupla do anel: trissubstituída, mais rica em elétrons.
\textbf{21.} 1,2-Epóxi-1-metil-4-(prop-1-en-2-il)ciclo-hexano, com o oxigênio em
ponte entre C1 e C2.
\textbf{22.} Em C1, o carbono mais substituído (que leva a metila).
\textbf{23.} 1-Metil-4-(prop-1-en-2-il)ciclo-hexano-1,2-diol, com os dois \ce{OH} em
trans (abertura anti).
\textbf{24.} Uma \omterm{def:b2:alkene-redox:dihydroxylation}{di-hidroxilação} syn da ligação dupla do anel (\ce{OsO4}
catalítico), que dá o diol cis --- desde que se consiga fazer o reagente
reagir seletivamente com essa ligação dupla.
\textbf{25.} \qty{1.00}{g} de limoneno absorve $\boldsymbol{\approx \qty{364}{mL}}$ de
di-hidrogênio.
\end{solution}
